\documentclass[11pt,a4paper]{article}

\usepackage[
    top= 1in,
    right=0.5in,
    bottom=1in,
    left=0.5in]{geometry}
\usepackage[utf8]{inputenc}
\usepackage[english]{babel}
\usepackage[T1]{fontenc}
\usepackage{lmodern}
\usepackage{minted}
\usepackage{amsmath,amssymb,amsfonts}
\usepackage{graphicx}

\title{Extreme Multiprogramming\\Assignment 2 - programming}
\author{Malte Stær Nissen}

\begin{document}
\date{\today}

\maketitle
\section{Graphical CSP-network design}
Figure \ref{fig:model} shows a CSP network graph of the MUD game implemented in
this assignment. When the graph says $\text{player}[x]$ it refers to a
specific player $x$ and the information/channels associated with this player in
the specific room. The game can be run by executing the file \texttt{game.py}
with python 2.7.4 and PyCSP installed.
\begin{figure}
    \centering
    \includegraphics[width=0.8\textwidth]{network.pdf}
    \caption{MUD CSP network graph/model}
    \label{fig:model}
\end{figure}
\section{Considerations on concurrency in the game}
The system is set up to let the room processes perform all data-handling.
Each client has a corresponding user object with a reading end of a channel
for reading user commands and a writing end for another channel returning
data/control to the client once the room has processed the given command.
Furthermore each client has a notification channel and a corresponding
notification reader process which receives notifications of other clients
entering and leaving the room which the client is currently in. The
notification retrieval is performed in an asynchronous process to avoid
deadlocks when notifying people in other rooms.

Each room has an ``input'' channel (reading end), which all the adjacent
rooms can write to in order to transfer a client into the room. The actual
transfer (sender end) is performed in an asynchronous process in order to
avoid deadlocks if two rooms try to transfer players to one another at the
same time. A transfer is performed by first removing the player ``object''
from the room it resides and then write it to the input channel of the new
room followed by notifications to all players in question and a final return
signal (True) to the client return channel.

The clients are purposely set to sleep for 3 seconds in between walking commands
in order to avoid cluttering std.out with notifications. This can be removed in
\texttt{agent.py}.

\section{Shutdown of the system}
Since each room has a list of clients and client notification processes as
well as channels to the adjacent rooms, a shutdown is a simple poison of all
channels the room is associated with. This will eventually
force all rooms to poison all their clients and thus shutting down the entire
network as long as all rooms were (in)directly connected at startup (which they
are in my initialization).


\clearpage
\appendix

\section{room.py}
\inputminted[linenos,fontsize=\scriptsize]{python}{src/room.py}

\section{agent.py}
\inputminted[linenos,fontsize=\scriptsize]{python}{src/agent.py}

\section{player.py}
\inputminted[linenos,fontsize=\scriptsize]{python}{src/player.py}

\section{game.py}
\inputminted[linenos,fontsize=\scriptsize]{python}{src/game.py}

\end{document}

